\documentclass{article}
\usepackage[T1]{fontenc}      % Meilleure gestion des accents en PDF
\usepackage[utf8]{inputenc}    % Encodage UTF-8 (par sécurité)
\usepackage{lmodern}           % Police moderne (bonne qualité)
\usepackage{amsmath, amssymb}  % Maths
\usepackage{mathtools}         % Extensions utiles (meilleur alignement)
\usepackage{enumitem}          % Listes personnalisées
\usepackage{geometry}          % Marges propres
\geometry{margin=2.5cm}


\title{Matrices}
\author{David Thébault}
\date{\today}

\begin{document}
\maketitle

\section{Produits de matrice}
\subsection{Définition}
Une matrice $A$ de taille $m\times n$ est un tableau de $m$ lignes et $n$ colonnes, dont les éléments sont des nombres réels (ou complexes). On note $A = (a_{ij})$ où $a_{ij}$ est l'élément situé à la ligne $i$ et à la colonne $j$.
\subsection{Systèmes d'équations linéaires}
Les matrices apparaissent dasn l'étude des systèmes d'équations linéaires. Un système d'équations linéaires peut être représenté sous la forme matricielle $AX = b$, où $A$ est la matrice des coefficients, $X$ est le vecteur des inconnues, et $b$ est le vecteur des constantes.

Par exemple, le système d'équations suivant :
\[
\begin{cases}
b_1 = x_1 - 2x_2 + 3x_3, \\
b_2 = -2x_1 + 4x_2 - 6x_3
\end{cases}
\]
peut être représenté par la matrice 
\[
A = \begin{pmatrix}1 & -2 & 3 \\ -2 & 4 & -6\end{pmatrix} \in \mathbb{R}^{2\times 3}, \qquad 
X = \begin{pmatrix}x_1 \\ x_2 \\ x_3\end{pmatrix} \in \mathbb{R}^{3}, \qquad
b = \begin{pmatrix}b_1 \\ b_2\end{pmatrix} \in \mathbb{R}^{2}.
\]

À la matrice $A$ est associée une application linéaire 
\[
f : \mathbb{R}^3 \to \mathbb{R}^2, \qquad 
x \mapsto f(x) = 
\begin{pmatrix}
x_1 - 2x_2 + 3x_3 \\
-2x_1 + 4x_2 - 6x_3
\end{pmatrix}.
\]

Ainsi résoudre le système $AX = B$ revient donc à trouver les antécédents de $b$ par l'application linéaire $f$ associée à la matrice $A$.
\[
f^{-1}(\{b\}) := \{x \in \mathbb{R}^3 \mid f(x) = b\}
\]
Image réciproque de $b$ par $f$. Cette image est non vide si et seulement si le système est compatible i.e. admet au moins une solution.
\[
f(\mathbb{R}^3) := \{f(x) \mid x \in \mathbb{R}^3\} \subset \mathbb{R}^2
\]
\section{Conclusion}
\end{document}